% GENERATED FILE. Edit canonical lessons, then run npm run generate:books.
% canonical-source-hash: fnv1a64:271b52a6bda925e0
% canonical-lessons: ZH-C76-yihuir, ZH-W76-hui, ZH-C76-mingtianjian, ZH-C76-zaishuo, ZH-C76-wanan, ZH-C76-baibai

\chapter{Parting}
\label{ch:zh-parting}
\hlchaptermodality{76}

\begin{chapteropening}
\textbf{By the end of this chapter:} \emph{I can say see you soon, see you tomorrow, good night and bye-bye.}

Complete the last lesson of chapter 76: I can say see you soon, see you tomorrow, good night and bye-bye.
\end{chapteropening}

\section[\texorpdfstring{\zh{一会儿} (yíhuìr)}{一会儿 (yíhuìr)}]{\zh{一会儿} (yíhuìr) — a little while; in a moment}

\label{lesson:ZH-C76-yihuir}



\begin{quote}
Before the new one: say the Mandarin for present, then the Mandarin for to give as a present.
\end{quote}

\subsection*{You'll want to know: \zh{一会儿}}
\textbf{\zh{一会儿}} — \emph{yíhuìr} — \textquotedblleft{}a little while; in a moment\textquotedblright{}.

\begin{quote}
\textbf{\zh{我一会儿打电话}} \emph{wǒ yíhuìr dǎ diànhuà} — I'll ring in a little while.
\end{quote}

The new character, in pieces: \zh{会} — \zh{人} spread as a roof, two bars, and a bend with a dot. The next lesson puts it in your hand.

yī changes its tone here, rising to yí before the falling huì: yíhuìr.

\begin{cousinweb}
\textbf{\zh{一会儿}} (\emph{yíhuìr}) is built from \zh{一} 'one', huì and the ending r, inherited from Middle Chinese, \emph{'jit}, \emph{hwajH} and \emph{nye} (Guangyun, in Baxter's transcription). The dictionary glosses huì as 'to come together'; how the phrase came to mean a short while is not settled.
\end{cousinweb}

\begin{grammarlens}[title={What you've built}]
One more word for times and partings.
\end{grammarlens}

\subsection*{Your turn}
\begin{itemize}
\raggedright
  \item \emph{Say it:} \emph{yíhuìr}
  \item \emph{Say it:} \emph{yíhuìr}, once more
  \item \emph{Say it:} \emph{wǒ yíhuìr dǎ diànhuà} — I'll ring in a little while.
  \item \emph{Recall:} say the Mandarin for present, then the Mandarin for to give as a present, then say \emph{yíhuìr} again
\end{itemize}

\begin{tcolorbox}[breakable,colback=teal!4,colframe=teal!35!black,title={Before you move on}]
What does \zh{一会儿} mean? (\textquotedblleft{}A little while; in a moment\textquotedblright{}.) Say it once more.
\end{tcolorbox}
\section[\texorpdfstring{\zh{会} (huì)}{会 (huì)}]{\zh{会} — \zh{人} over two bars and a bend}

\label{lesson:ZH-W76-hui}



\begin{quote}
The word first, then the character inside it.

\begin{itemize}
\raggedright
  \item \emph{Recall:} say \emph{yíhuìr} — \textbf{R1}, one lesson back, before it decays
\end{itemize}

\begin{quote}
\zh{人} \zh{今} \zh{去}
\end{quote}

\emph{Read it:} each one aloud, with nothing said first — \emph{rén}, \emph{jīn}, \emph{qù}
\end{quote}

\subsection*{Script — \zh{人} over two bars and a bend}
\hlblockfigure{\detokenize{figures/ZH-W76-hui-filmstrip.pdf}}{How \zh{会} is written, stroke by stroke}

\begin{quote}
\zh{会}
\end{quote}

\textbf{Three parts, top to bottom.}

\textbf{\zh{人}, two strokes}, spread wide as a roof, as on \zh{今}.

\textbf{Two horizontals} under it, the lower one longer.

\textbf{A stroke that falls left and turns right} along the bottom, then \textbf{a dot}: the same two strokes that end \zh{去}.

2 + 2 + 2 = six.

\subsection*{Your turn}
\begin{itemize}
\raggedright
  \item \emph{Write it:} \zh{会}, all six, with the model in front of you
  \item \emph{Cover:} write it again from memory, and count the pieces, not the strokes
\end{itemize}

\begin{tcolorbox}[breakable,colback=teal!4,colframe=teal!35!black,title={Before you move on}]
How many strokes in \zh{会}? (\textbf{Six}.) Which character you know is on top? (\textbf{\zh{人}}.) What comes last? (\textbf{A dot}.)

Source: \href{https://raw.githubusercontent.com/chanind/hanzi-writer-data/68d10a4b21150cae5e1ebbd223eed289cf32d90c/data/\%E4\%BC\%9A.json}{Hanzi Writer Data \zh{会}, pinned stroke paths and medians}.
\end{tcolorbox}
\section[\texorpdfstring{\zh{明天见} (míngtiān jiàn)}{明天见 (míngtiān jiàn)}]{\zh{明天见} (míngtiān jiàn) — see you tomorrow}

\label{lesson:ZH-C76-mingtianjian}



\begin{quote}
Before the new one: say the Mandarin for to give as a present, then the Mandarin for a little while.
\end{quote}

\subsection*{You'll want to know: \zh{明天见}}
\textbf{\zh{明天见}} — \emph{míngtiān jiàn} — \textquotedblleft{}see you tomorrow\textquotedblright{}.

\begin{quote}
\textbf{\zh{下雨我们明天见}} \emph{xiàyǔ wǒmen míngtiān jiàn} — It's raining: let's meet tomorrow.
\end{quote}

\begin{cousinweb}
\textbf{\zh{明天见}} (\emph{míngtiān jiàn}) is built from \zh{明天} 'tomorrow' and \zh{见} 'to see', the \zh{见} of \zh{再见}.
\end{cousinweb}

\begin{grammarlens}[title={What you've built}]
One more word for times and partings.
\end{grammarlens}

\subsection*{Your turn}
\begin{itemize}
\raggedright
  \item \emph{Say it:} \emph{míngtiān jiàn}
  \item \emph{Say it:} \emph{míngtiān jiàn}, once more
  \item \emph{Say it:} \emph{xiàyǔ wǒmen míngtiān jiàn} — It's raining: let's meet tomorrow.
  \item \emph{Recall:} say the Mandarin for to give as a present, then the Mandarin for a little while, then say \emph{míngtiān jiàn} again
\end{itemize}

\begin{tcolorbox}[breakable,colback=teal!4,colframe=teal!35!black,title={Before you move on}]
What does \zh{明天见} mean? (\textquotedblleft{}See you tomorrow\textquotedblright{}.) Say it once more.
\end{tcolorbox}
\section[\texorpdfstring{\zh{再说} (zàishuō)}{再说 (zàishuō)}]{\zh{再说} (zàishuō) — to leave it till later}

\label{lesson:ZH-C76-zaishuo}



\begin{quote}
Before the new one: say the Mandarin for a little while, then the Mandarin for see you tomorrow.
\end{quote}

\subsection*{You'll want to know: \zh{再说}}
\textbf{\zh{再说}} — \emph{zàishuō} — \textquotedblleft{}to leave it till later\textquotedblright{}.

\begin{quote}
\textbf{\zh{对不起一会儿再说}} \emph{duìbuqǐ yíhuìr zàishuō} — Sorry: let's talk about it in a bit.
\end{quote}

\begin{cousinweb}
\textbf{\zh{再说}} (\emph{zàishuō}) is built from \zh{再} 'again' and \zh{说} 'to say', both inherited from Middle Chinese, \emph{tsojH} and \emph{sywet} (Guangyun, in Baxter's transcription): say it again, another time.
\end{cousinweb}

\begin{grammarlens}[title={What you've built}]
One more word for times and partings.
\end{grammarlens}

\subsection*{Your turn}
\begin{itemize}
\raggedright
  \item \emph{Say it:} \emph{zàishuō}
  \item \emph{Say it:} \emph{zàishuō}, once more
  \item \emph{Say it:} \emph{duìbuqǐ yíhuìr zàishuō} — Sorry: let's talk about it in a bit.
  \item \emph{Recall:} say the Mandarin for a little while, then the Mandarin for see you tomorrow, then say \emph{zàishuō} again
\end{itemize}

\begin{tcolorbox}[breakable,colback=teal!4,colframe=teal!35!black,title={Before you move on}]
What does \zh{再说} mean? (\textquotedblleft{}To leave it till later\textquotedblright{}.) Say it once more.
\end{tcolorbox}
\section[wǎn'ān]{wǎn'ān — good night}

\label{lesson:ZH-C76-wanan}



\begin{quote}
Before the new one: say the Mandarin for see you tomorrow, then the Mandarin for to leave it till later.
\end{quote}

\subsection*{You'll want to know: wǎn'ān}
\emph{wǎn'ān} — \textquotedblleft{}good night\textquotedblright{}.

\begin{quote}
\emph{māma wǎn'ān} — Mum, good night.
\end{quote}

Its characters come later; for now it is a word for the ear.

\begin{cousinweb}
\emph{wǎn'ān} is built from wǎn 'evening' and ān 'peace', both inherited from Middle Chinese, \emph{mjonX} and \emph{'an} (Guangyun, in Baxter's transcription); the dictionary glosses the second as 'calm, at rest'.
\end{cousinweb}

\begin{grammarlens}[title={What you've built}]
One more word for times and partings.
\end{grammarlens}

\subsection*{Your turn}
\begin{itemize}
\raggedright
  \item \emph{Say it:} \emph{wǎn'ān}
  \item \emph{Say it:} \emph{wǎn'ān}, once more
  \item \emph{Say it:} \emph{māma wǎn'ān} — Mum, good night.
  \item \emph{Recall:} say the Mandarin for see you tomorrow, then the Mandarin for to leave it till later, then say \emph{wǎn'ān} again
\end{itemize}

\begin{tcolorbox}[breakable,colback=teal!4,colframe=teal!35!black,title={Before you move on}]
What does wǎn'ān mean? (\textquotedblleft{}Good night\textquotedblright{}.) Say it once more.
\end{tcolorbox}
\section[bàibai]{bàibai — bye-bye}

\label{lesson:ZH-C76-baibai}



\begin{quote}
Before the new one: say the Mandarin for to leave it till later, then the Mandarin for good night.
\end{quote}

\subsection*{You'll want to know: bàibai}
\emph{bàibai} — \textquotedblleft{}bye-bye\textquotedblright{}.

\begin{quote}
\emph{bàibai míngtiān jiàn} — Bye, see you tomorrow!
\end{quote}

Its characters come later; for now it is a word for the ear.

\begin{cousinweb}
\emph{bàibai} is borrowed from English \emph{bye-bye}, written with a character read bài, 'to bow', for the sound.
\end{cousinweb}

\begin{grammarlens}[title={What you've built}]
One more word for times and partings.
\end{grammarlens}

\subsection*{Your turn}
\begin{itemize}
\raggedright
  \item \emph{Say it:} \emph{bàibai}
  \item \emph{Say it:} \emph{bàibai}, once more
  \item \emph{Say it:} \emph{bàibai míngtiān jiàn} — Bye, see you tomorrow!
  \item \emph{Recall:} say the Mandarin for to leave it till later, then the Mandarin for good night, then say \emph{bàibai} again
\end{itemize}

\begin{tcolorbox}[breakable,colback=teal!4,colframe=teal!35!black,title={Before you move on}]
What does bàibai mean? (\textquotedblleft{}Bye-bye\textquotedblright{}.) Say it once more.
\end{tcolorbox}
